\begin{examplebox}
	\textbf{\textcolor{CExample}{例 1（基础）}}

	\textbf{\textcolor{CExample}{题目：}} 在直角三角形$\triangle ABC$中，$\angle C = 90^\circ$，斜边$AB=10$，$M$是$AB$的中点。求线段$CM$的长。

	\textbf{\textcolor{CExample}{解题步骤：}}
	\begin{description}[style=nextline, leftmargin=2.8em, labelsep=0.5em, itemsep=0.45em]
		\item[\textbf{第一步（识别条件）：}] 检查是否满足定理条件：直角三角形，$M$是斜边中点。
		      \par\textbf{理由：}
		      定理要求$\angle C=90^\circ$且$M$为斜边$AB$的中点，题目满足。

		\item[\textbf{第二步（套用定理）：}] 由直角三角形斜边中线定理，
		      \[
			      CM = \frac12 AB.
		      \]
		      \par\textbf{理由：}
		      直角三角形斜边上的中线等于斜边的一半。

		\item[\textbf{第三步（计算）：}] 代入$AB=10$，
		      \[
			      CM = \frac12 \times 10 = 5.
		      \]
		      \par\textbf{理由：}
		      直接乘法计算。
	\end{description}

	\textbf{\textcolor{CExample}{关键结论：}} $CM=5$。

	\medskip
	\textbf{\textcolor{CExample}{例 2（稍变形）}}

	\textbf{\textcolor{CExample}{题目：}} 在直角三角形$\triangle ABC$中，$\angle C = 90^\circ$，$CM$为斜边$AB$上的中线，且$CM=6$。求斜边$AB$的长。

	\textbf{\textcolor{CExample}{解题步骤：}}
	\begin{description}[style=nextline, leftmargin=2.8em, labelsep=0.5em, itemsep=0.45em]
		\item[\textbf{第一步（识别条件）：}] 确认$CM$是斜边$AB$上的中线。
		      \par\textbf{理由：}
		      题目已给出$\angle C=90^\circ$，且$CM$为斜边$AB$上的中线，满足定理条件。

		\item[\textbf{第二步（正向定理的等式变形）：}] 由直角三角形斜边中线定理，
		      \[
			      CM = \frac12 AB.
		      \]
		      两边同乘$2$，得
		      \[
			      AB = 2CM.
		      \]
		      \par\textbf{理由：}
		      这里不是使用逆定理，而是把正向定理的等式进行变形。

		\item[\textbf{第三步（计算）：}] 代入$CM=6$，
		      \[
			      AB = 2 \times 6 = 12.
		      \]
		      \par\textbf{理由：}
		      乘法计算。
	\end{description}

	\textbf{\textcolor{CExample}{关键结论：}} $AB=12$。

	\medskip
	\textbf{\textcolor{CExample}{例 3（综合应用）}}

	\textbf{\textcolor{CExample}{题目：}} 在直角三角形$\triangle ABC$中，$\angle C=90^\circ$，$\angle A=30^\circ$，斜边$AB=10$，$M$是$AB$的中点。求$CM$的长以及$\triangle ABC$的周长和面积。

	\textbf{\textcolor{CExample}{解题步骤：}}
	\begin{description}[style=nextline, leftmargin=2.8em, labelsep=0.5em, itemsep=0.45em]
		\item[\textbf{第一步（求中线长）：}] 由直角三角形斜边中线定理，
		      \[
			      CM = \frac12 AB = \frac12 \times 10 = 5.
		      \]
		      \par\textbf{理由：}
		      斜边中线等于斜边的一半。

		\item[\textbf{第二步（求直角边）：}] 在直角三角形$\triangle ABC$中，$\angle A=30^\circ$，所以
		      \[
			      BC = AB \cdot \sin30^\circ = 10 \times \frac12 = 5,
		      \]
		      \[
			      AC = AB \cdot \cos30^\circ = 10 \times \frac{\sqrt3}{2} = 5\sqrt3.
		      \]
		      \par\textbf{理由：}
		      $BC$是$\angle A$的对边，$AC$是$\angle A$的邻边，按直角三角形锐角三角函数定义计算。

		\item[\textbf{第三步（求周长和面积）：}]
		      周长和面积计算如下：
		      \[
			      \begin{aligned}
				      C_{\triangle} & = AB + BC + AC     \\
				                    & = 10 + 5 + 5\sqrt3 \\
				                    & = 15 + 5\sqrt3,
			      \end{aligned}
		      \]
		      \[
			      \begin{aligned}
				      S_{\triangle} & = \frac12 \times BC \times AC     \\
				                    & = \frac12 \times 5 \times 5\sqrt3 \\
				                    & = \frac{25\sqrt3}{2}.
			      \end{aligned}
		      \]
		      \par\textbf{理由：}
		      周长等于三边之和；直角三角形面积等于两直角边乘积的一半。
	\end{description}

	\textbf{\textcolor{CExample}{关键结论：}} $CM=5$，周长为$15+5\sqrt3$，面积为$\frac{25\sqrt3}{2}$。
\end{examplebox}